\BourbakiBackSection{HISTORICAL NOTE}

(Numbers in brackets refer to the bibliography at the end of this note.)

Every measurement of quantities implies a vague notion of real numbers (we shall see the exact reasons for this in Chapter V, § 2). From the mathematical point of view, the origins of the theory of real numbers can be traced back to the progressive formation by the Babylonians of a system of numeration which was (in principle) capable of representing arbitrarily close approximations to any real number {[}1{]}. The possession of such a system, and the confidence in numerical calculation which naturally resulted from it, inevitably led to a ``naïve'' notion of real number which differs hardly at all from that which is current today (linked with the decimal system of numeration) in elementary education and among physicists and engineers. This notion cannot be precisely defined, but can be expressed by saying that a number is regarded as defined by the possibility of finding approximations to it and using these approximations in calculation; this necessarily implies a certain amount of confusion between measures of physical quantities, which of course are not susceptible to an infinite series of successively closer and closer approximation, and ``numbers'' such as \(\sqrt{2}\) (assuming that one is in possession of an algorithm which would make possible an infinite series of successively closer and closer approximation of such numbers).

A similar ``pragmatic'' attitude appears in all mathematical schools in which expertise in calculation is more important than rigour and theory. The latter, however, were predominant in Greek mathematics; and it is to the Greeks that we owe the first rigorous and coherent theory of ratios of magnitudes, that is, essentially, of real numbers. This theory was the culmination of a series of discoveries about proportions and, in particular, incommensurable ratios, whose importance in the history of Greek thought can hardly be exaggerated, but which in the absence of accurate texts can be discerned only in outline. Greek mathematics in its early stages was inextricably bound up with speculations, part scientific and part philosophical and mystical, about proportion, similitude and ratio, especially ``simple ratios'' (expressible by fractions with small numerators and denominators); and one of the characteristic tendencies of the Pythagorean school was to attempt to explain all in terms of integers and ratios of integers. But it was the Pythagorean school, in fact, which discovered that the diagonal of a square is incommensurable with its side (in other words, that \(\sqrt{2}\) is irrational). This is without doubt the first example of a proof of impossibility in mathematics, and the mere fact of posing such a question implies a clear distinction between a ratio and approximations to it, and indicates the immense gap which separates the Greek mathematicians from their predecessors (*).

We know little about the movement of ideas which accompanied and followed this important discovery ( \(^{**}\) ). We shall give only a brief summary of the main ideas which lie at the base of the theory of ratios of magnitudes, which was constructed by the great mathematician Eudoxus (a contemporary and friend of Plato), definitively adopted by classical Greek mathematics, and is known to us through Euclid's Elements {[}2{]} where it is described in a masterly fashion (in Book V of the Elements):

\begin{enumerate}
\def\labelenumi{\arabic{enumi})}
\item
  The word and the idea of number are strictly reserved to natural integers \textgreater1 (1 is the monad and not, strictly speaking, a number), to the exclusion not only of our irrational numbers but also of what we call rational numbers: to the Greek mathematicians of the classical period the latter are ratios of numbers. There is much more here than a simple question of terminology: the word ``number'' was for the Greeks (and for the moderns up to a recent time) linked with the idea of a system with two laws of composition (addition and multiplication); ratios of integers were regarded by the classical Greek mathematicians as operators, defined on the set of integers or on some subset of this set {[}the ratio of p to q is the operator which, applied to N, if N is a multiple of q, gives the integer \(p.(N/q)\) {]}, and forming a multiplicative group but not a system with two laws of composition. In this the Greek mathematicians separated themselves voluntarily from the ``logisticians'' or professional calculators who, like their Egyptian and Babylonian predecessors, had no scruples about treating fractions as if they were numbers, or adding a fraction to an integer. It seems moreover that this self-imposed restriction on the concept of number came from philosophical rather than mathematical motives, and followed the reflections of the first Greek thinkers on the unit and the multiple; the unit (in this system of thought) being incapable of subdivision without thereby losing its character of unit (*).
\item
  The theory of magnitudes is based on axioms, which applied simultaneously to all types of magnitudes (there are allusions to earlier theories which apparently treated lengths, areas, volumes, times, etc., all separately). Magnitudes of the same type are characterized by the facts that they can be compared (that is to say, it is assumed that equality, which is an equivalence relation, and the relations \textgreater{} and \textless{} are defined), that they can be added and subtracted (A + B is defined, and so is A - B if A \textgreater{} B) and that they satisfy ``Archimedes' axiom'' (Theorem 1 of § 2). It is clearly realized from the beginning that this latter fact is the keystone of the whole edifice (it is in fact indispensable in any axiomatic characterization of real numbers; cf.~Chapter V, § 2). Its attribution to Archimedes is purely accidental: in the introduction to his ``Quadrature of the Parabola'' {[}3{]}, Archimedes emphasizes that this axiom was used by his predecessors, that it played an essential part in the work of Eudoxus, and that its consequences were no less certainly established than determinations of areas and volumes performed without its help (**).
\end{enumerate}

We shall see in Chapter V, § 2, how the theory of real numbers is a necessary consequence of this axiomatic foundation. Note that, for Eudoxus, the magnitudes of a given type form a system with one internal law of composition (addition), but that this system has an external law of composition whose operators are ratios of magnitudes, conceived of as forming an abelian multiplicative group. If A and A' are magnitudes of the same type, and if B and B' are magnitudes of the same type, then the ratios of A to A' and B to B' are defined to be equal if, for all integers m and m', the relation mA \textless{} m'A' implies mB \textless{} m'B', and mA \textgreater{} m'A' implies mB \textgreater{} m'B'. Inequalities between ratios are defined by similar methods. That these ratios form a domain of operators for every type of magnitude is equivalent to the axiom (not explicitly stated but frequently used in Euclid's exposition) of the existence of the fourth proportional: if a ratio A/A' is given, and if B is given, then there exists a B', of the same type as B, such that B/B' = A/A'. Thus Eudoxus' brilliant idea allowed him to identify all the domains of operators defined by different types of magnitude (*); in a similar way, the set of ratios of integers (see above) can be identified with a subset of the set of ratios of magnitudes, namely the set of rational ratios (or ratios of commensurable magnitudes); nevertheless, since these ratios, regarded as operators of the integers, are (in general) defined only on a subset of the set of integers, it was necessary to develop their theory separately (Book VII of Euclid).

The universal domain of operators thus constructed was the equivalent, for the Greek mathematicians, of what the set of real numbers is for us; moreover it is clear that, with addition of magnitudes and multiplication of ratios of magnitudes, they possessed the equivalent of what the field of real numbers is for us, although in a much less manageable form ( \(^{**}\) ). On the other hand, one may ask whether they regarded these sets (the set of magnitudes of a given type, or the set of ratios of magnitudes) as complete in our sense. It is not clear, otherwise, why they should have assumed the existence of the fourth proportional (without even perceiving the need to make it an axiom); moreover, some texts seem to refer to ideas of this nature; and they certainly assumed as self-evident the fact that a curve which can be described by continuous motion cannot pass from one side of a line to the other without cutting the line, a principle which they used, for example, in their investigations on the duplication of the cube (construction of \(\sqrt[3]{2}\) by intersections of curves) and which is essentially equivalent to the property in question; nevertheless, the texts which have come down to us do not enable us to discern their ideas on this point with complete accuracy.

Such was the state of the theory of real numbers in the classical period of Greek mathematics. Admirable though Eudoxus' construction was, and leaving nothing to be desired in rigour or coherence, nevertheless it must be admitted that it lacked flexibility and did not encourage the development of numerical calculation, still less the development of algebraic calculation. Moreover its logical necessity could not be apparent except to those in love with rigour and familiar with abstraction; thus it is natural that, with the decline of Greek mathematics, the ``naïve'' point of view, which had been preserved through the tradition of the logisticians, should gradually re-emerge. This point of view is dominant for example in Diophantus {[}4{]}, who in truth was an upholder of this tradition rather than of official Greek science; he perfunctorily reproduces the Euclidean definition of number, but in reality he uses the word ``number'' to mean the unknown in algebraic problems whose solution may be either an integer, or a fraction, or an irrational number (*). Although this change of attitude on the subject of number is connected with one of the most important advances in the history of mathematics, namely the development of algebra, it does not of course constitute an advance in itself, but rather a retreat.

We cannot trace here the vicissitudes of the concept of number through Hindu, Arab and western mathematics up to the end of the Middle Ages. The ``naïve'' notion of number predominated, and although the Elements of Euclid served as a basis for the teaching of mathematics during this period, it is most likely that the doctrine of Eudoxus remained generally uncomprehended because the need for it was no longer appreciated. The ``ratios'' of Euclid were customarily described as ``numbers'', and the rules for calculating with integers were applied to them without any attempt to analyse the reasons for the success of these methods.

Nevertheless we see R. Bombelli, as early as the middle of the 16th century, expounding a point of view on this subject, in his Algebra {[}5{]} (*), which is essentially correct (provided that the results of Book V of Euclid are assumed to be known); having realized that once the unit of length has been chosen there is a one-to-one correspondence between lengths and ratios of magnitudes, he defines the various algebraic operations on lengths (assuming of course that the unit has been fixed) and, representing numbers by lengths, obtains the geometrical definition of the field of real numbers (a point of view which is usually credited to Descartes) and thus gives his algebra a solid geometrical foundation (**).

But Bombelli's Algebra, though singularly advanced for its time, did not go beyond the extraction of radicals and the solution by radicals of equations of the second, third and fourth degrees; of course the possibility of extraction of radicals is assumed without any discussion. Simon Stevin {[}6{]} adopts an analogous viewpoint on the subject of numbers; for him a number denotes a measure of a magnitude and is regarded as being essentially ``continuous'' (without giving a precise meaning to this word). If he distinguishes between ``geometric numbers'' and ``arithmetic numbers'', it is only because of the accident of their mode of definition and not because of a difference in nature. His last words on this subject run as follows: ``We conclude therefore that there are no absurd, irrational, irregular, inexplicable or surd numbers, but that there is in them such excellence and concordance that we have matter for meditation night and day on their admirable perfection'' ({[}6{]}, p.~10). On the other hand he was the first to use decimal fractions as a method of calculation and proposed a notation for them which is very close to ours; and he saw clearly that these fractions provide an algorithm for indefinitely close approximation to any real number, as is shown by his Appendice algebraique of 1594, ``containing a general rule for all equations'' (the only known copy of this pamphlet was burnt at Louvain in 1914; but see {[}6{]}, vol.~1, p.~88). Such an equation being written in the form \(\mathrm{P}(x) = \mathrm{Q}(x)\) {[}where the degree of the polynomial \(\mathbf{P}\) is greater than the degree of the polynomial \(\mathbf{Q}\) , and \(\mathrm{P(0)} < \mathrm{Q(0)}\) {]}, he substitutes for \(x\) the numbers 10, 100, 1000, \ldots{} until he finds \(\mathrm{P}(x) > \mathrm{Q}(x)\) which, he says, determines the number of digits in the root; then (if for example the root has two digits) he substitutes 10, 20, \ldots{} which determines the number of tens; then similarly for the number of units, then for the successive decimal digits: ``And proceeding indefinitely in this way'', he says, ``we approach infinitely near to the number required'' ({[}6{]}, p.~88). As we see, Stevin had exactly the idea of Theorem 2 of § 6 (and was without doubt the first to have it) and recognized that this theorem was the essential tool for the systematic solution of numerical equations; at the same time we can see that he had so clear an intuitive conception of the numerical continuum that little remained to be done to make it definitively precise.

Nevertheless, in the following two centuries the definitive establishment of correct methods was twice retarded by the development of two theories whose history does not come within the framework of this note: the infinitesimal calculus and the theory of series. Throughout the discussions they gave rise to, we perceive, as in all periods of the history of mathematics, the perpetual balance between those who sought to push forward, at the cost of some insecurity, convinced that there would always be time later to consolidate the conquered terrain, and those of a critical mind who (without necessarily being in any way inferior in point of intuitive gifts and inventive genius) believed that their energies were not wasted when they devoted their effort to the precise formulation and rigorous justification of their conceptions. In the seventeenth century the main subject of debate was the notion of ``infinitely'' small which, though justified a posteriori by the results which were obtained with its help, seemed to be in open opposition to the axiom of Archimedes; and we see the most enlightened minds of this period finally adopting a point of view which differed little from that of Bombelli, and which is distinguished above all by the greater attention it paid to the rigorous methods of the ancients. Isaac Barrow (Newton's teacher, who himself played an important part in the creation of the infinitesimal calculus) gave a brilliant exposition of this viewpoint in his Mathematical Lectures given at Cambridge in 1664-5-6 {[}7{]}; recognizing the need to return to the theory of Eudoxus in order to regain the proverbial ``geometrical certainty'' in the subject of number, he presents at length an extremely judicious defence of Eudoxus' theory (which, on Barrow's evidence, had remained unintelligible to many of his contemporaries) against those who charged it with obscurity or even absurdity. On the other hand, defining numbers to be symbols which denote ratios of magnitudes and to be capable of being combined by the operations of arithmetic, Barrow obtains the field of real numbers in terms which Newton took up again in his Arithmetic, and which his successors up to Dedekind and Cantor did not change.

But it was in this period that the method of expansion in series was introduced; this rapidly took on an exclusively formal character in the hands of impenitent algebraists and deflected the attention of mathematicians from the questions of convergence which are essential to any sound use of series in the domain of real numbers. Newton, the principal creator of the method, was certainly aware of the necessity of considering these questions; and if he did not elucidate them sufficiently, he at least realized that the power series which he introduced converged ``usually'' at least as well as a geometric series (whose convergence was already known to the ancients) for small values of the variable {[}8{]}. At about the same time, Leibniz had observed that an alternating series whose terms decrease in absolute value and tend to 0 is convergent. In the following century, d'Alembert in 1768 raised doubts about the use of non-convergent series; but the authority of the Bernoullis and above all Euler was such as to make doubts of this nature exceptional at this period.

It is clear that mathematicians who were in the habit of using series for numerical calculations would never have so neglected the notion of convergence; and it is no accident that the first to lead the return to correct methods, in this domain as in many others, was a mathematician who in his early youth was in love with numerical calculation. C.F. Gauss, who, when still almost a child, had practised the algorithm of the arithmeticogeometric mean (*), could scarcely fail to form a clear notion of the limit; and in a fragment dated 1800 (but first published in our times) ({[}9{]}, Vol. X¹, p.~390) he gives precise definitions of the least upper bound and the greatest lower bound, and the upper and lower limits of a sequence of real numbers; the existence of the bounds (for a bounded sequence) he seems to have assumed as obvious, and the upper and lower limits he defines correctly as the limits, as n tends to +∞, of sup \(u_{n+p}\) and inf \(u_{n+p}\) respectively. On the other hand, in his memoir of 1812 on the hypergeometric series ({[}9{]}, vol.~III, p.~139), Gauss gives the first model of a discussion of convergence conducted, as he says, ``in full rigour, and made to satisfy those whose preferences are for the rigorous methods of the ancient geometers''. It is true that this discussion, which occupies a secondary place in the memoir, does not go back to the first principles of the theory of series; Cauchy, in his Cours d'Analyse of 1821 {[}10{]}, was the first to establish these in a manner correct in every point, starting from Cauchy's criterion (clearly stated, and assumed as self-evident). Since he takes the point of view of Barrow and Newton on the definition of number, it can be said that for Cauchy the real numbers are defined by the axioms of magnitudes and Cauchy's criterion; and this is in fact enough to define them (see Chapter V, § 2).

(*) If \(x_0, y_0\) are given and \(> 0\) , let \(x_{n+1} = (x_n + y_n)/2\) and

\[
y _ {n + 1} = \sqrt {x _ {n} y _ {n}};
\]

At the same moment another important aspect of the theory of real numbers was definitively cleared up. As we have said, it had always been assumed as geometrically obvious that two continuous curves cannot cross each other without intersecting --- a principle which (suitably made precise) is again equivalent to the completeness of the real line (as a uniform space). This principle is at the base of the ``rigorous'' proof, given by Gauss in 1799, of d'Alembert's theorem, according to which every polynomial with real coefficients has a real or complex root ({[}9{]}, vol.~III, p.~1); the proof of the same theorem given by Gauss in 1815 ({[}9{]}, vol.~III, p.~31) depends, as does an earlier attempt by Lagrange, on the analogous but simpler principle that a polynomial cannot change sign without vanishing; this is a particular case of Theorem 2 of § 6, which we have already seen used by Stevin. In 1817 Bolzano gave a complete proof, founded on Cauchy's criterion, of this latter principle, which he obtains as a particular case of the analogous theorem for continuous real-valued functions of a real variable {[}11{]}. He enunciates ``Cauchy's criterion'' clearly (and before Cauchy did), and seeks to justify it by an argument which, in the absence of any arithmetic definition of real numbers, was only and could only be a vicious circle; but, once this point has been accepted, his work is entirely correct and most remarkable, since it contains not only the modern definition of a continuous function (given here for the first time), with the proof of the continuity of polynomials, but also the proof of the existence of the greatest lower bound of an arbitrary bounded set of real numbers (he speaks not of sets, but of properties of real numbers, which comes to the same thing). Cauchy in his Cours d'Analyse {[}10{]} also defined continuous functions of one or more real variables and proved by the same reasoning as Simon Stevin used, that a continuous function of one variable cannot change sign without vanishing. This reasoning becomes correct, once continuity has been defined, if Cauchy's criterion is invoked (or if one assumes, as Cauchy does at this point, the equivalent principle of ``nested intervals''; the convergence of decimal fractions is of course only a particular case of this principle).

Once arrived at this point, there remained for the mathematicians only the task of making precise and extending the results obtained, by correcting various errors and filling various gaps. For example, Cauchy had believed at one time that a convergent series, whose terms are continuous functions of one variable, has a continuous function as its sum. Abel's rectification of this point, in the course of his important work on series ({[}12{]}, vol.~1, p.~219; cf.~also vol.~2, p.~257 and passim) led finally to the elucidation by Weierstrass of the concept of uniform convergence (in his lectures, which remained unpublished but which had a considerable influence; see the Historical Note to Chapter X). Again, Cauchy had assumed, without sufficient justification, the existence of the minimum of a continuous function in one of his proofs of the existence of roots of a polynomial; and again it was Weierstrass who threw light on questions of this nature by proving (in his lectures) Theorem 1 of § 6 for functions of real variables, defined on bounded closed intervals. Following his criticism of unjustified applications of this theorem to sets of functions (``Dirichlet's principle'' is the best-known example) there began the movement of ideas which led, as we have seen in the Historical Note to Chapter I, to the general definition of compact spaces and the modern statement of the theorem as we have given it.

At the same time Weierstrass, in his lectures, had perceived the logical importance in making the idea of real number entirely independent of the theory of magnitudes; the latter is effectively equivalent to an axiomatic definition of the points of the line (and thus of the set of real numbers) and the assumption of the existence of such a set; although this method is essentially correct, it is evidently preferable to start only from the rational numbers, and to construct the real numbers from them by completion (*). This was achieved, by diverse methods and independently of each other, by Weierstrass, Dedekind, Méray and Cantor; while the method of ``cuts'', proposed by Dedekind {[}13{]}, came very near to the definitions of Eudoxus, the other methods proposed are close to that which is expounded in this series. Simultaneously Cantor began to develop the theory of sets of real numbers, the idea of which was first conceived by Dedekind (see the Bibliography to Chapter I), and thus obtained the principal elementary results on the topology of the real line, the structure of its open and closed sets, the notion of derived set and of totally disconnected perfect set, etc.; he also obtained Theorem 1 of § 8 on the power of the continuum, and deduced from it that the continuum is uncountable, that the set of transcendental numbers has the power of the continuum, and also (a paradoxical result for its time) that the set of points of a plane (or of space) has the same power as the set of points of a line.

With Cantor, the questions studied in this chapter assumed practically their definitive form. We refer to the Historical Note on Chapter I for the immediate impact of Cantor's work; let us indicate briefly the directions in which it has extended. Apart from leading to work on general topology (see Chapter I) and applications to integration, which will be dealt with thoroughly elsewhere in this series, Cantor's work has led to investigations of the structure and classification of sets of points on a line, and of real-valued functions of a real variable. These have their origin in the work of Borel {[}14{]} which was directed mainly towards measure theory but which led, among other things, to the definition of ``Borel sets'', i.e., sets belonging to the smallest family of subsets of R which contains all intervals and is closed with respect to union and countable intersection and with respect to the operation C (cf.~Chapter IX, § 6, no. 3). These sets are closely related to the so-called ``Borel functions'' or ``Baire functions'', that is to say functions which can be obtained from continuous functions by the operation of taking the limit of a sequence, repeated ``transfinitely''; they were defined by Baire in the course of an important work in which he entirely abandoned the measure viewpoint for a systematic investigation of the qualitative and ``topological'' aspect of these questions {[}15{]}; in this context he defined and studied semi-continuous functions (and was the first to do so), and in order to characterize functions which are limits of continuous functions he introduced the important notion of a meagre set (set ``of the first category'' in Baire's terminology) which we shall study in Chapter IX. As to the many works which have followed Baire, and which are mainly products of the Russian and especially the Polish schools, we can do no more here than draw attention to their existence (see for example {[}16{]} and the journal Fundamentae Mathematicae).

\BourbakiBackSection{BIBLIOGRAPHY}

{[}1{]} O. NEUGEBAUER, Vorlesungen über Geschichte der antiken Mathematik, Bd. I: Vorgriechische Mathematik, Berlin (Springer), 1934.

{[}2{]} Euclidis Elementa, 5 volumes, ed.~J. L. Heiberg, Lipsiae (Teubner), 1883-88.

{[}2a{]} T. L. HEATH, The Thirteen Books of Euclid's Elements\ldots, 3 volumes, Cambridge, 1908.

{[}3{]} Archimedis Opera Omnia, 3 volumes, ed.~J. L. Heiberg, 2nd edition, 1913-15.

{[}3a{]} Les Œuvres complètes d'Archimède, translated by P. Ver Eecke, Paris-Bruxelles (Desclée-de Brouwer), 1921.

{[}4{]} Diophanti Alexandrini Opera Omnia\ldots, 2 volumes, ed.~P. Tannery, Lipsiae (Teubner), 1893-95.

{[}4a{]} Diophante d'Alexandrie, translated by P. Ver Eecke, Bruges (Desclée de Brouwer), 1926.

{[}5{]} R. BOMBELLI, L'Algebra, ed.~E. Bortolotti, Bologna (Zanichelli), 1929.

{[}6{]} LES ŒUVRES Mathématiques de SIMON STEVIN de Bruges, Ou sont insérées les MÉMOIRES MATHÉMATIQUES, Esquelles s'est exercé le Très-haut et Très-illustre Prince MAURICE DE NASSAU, Prince d'Aurenge, Gouverneur des Provinces des Païs-Bas unis, General par Mer et par Terre, etc., Le tout reveu, corrigé et augmenté par ALBERT GIRARD Samielois, Mathématicien, A LEYDE, Chez Bonaventure et Abraham Elsevier, Imprimeurs ordinaires de l'Université, Anno MDCXXXIV (= 1634), vol.~I.

{[}7{]} I. BARROW, Mathematical Works, Cambridge (University Press), 1860.

{[}8{]} I. NEWTON, De Analysi per equatione numero terminorum infinitas, in Commercium Epistolicum D. Johannis Collins et aliorum de Analysi promota, Londini, 1712.

{[}9{]} C. F. GAUSS, Werke, vols. III (Göttingen, 1816) and \(\mathbf{X}_1\) (ibid., 1917).

{[}10{]} A. CAUCHY, Cours d'Analyse de l'École Royale Polytechnique, 1re partie, 1821 = Œuvres (II) vol.~III, Paris (Gauthier-Villars), 1897.

{[}11{]} B. BOLZANO, Rein Analytischer Beweis des Lehrsatzes, dass zwischen je zwei Werthen, die ein entgegengesetztes Resultat gewähren, wenigstens eine reelle Wurzel liegt, Ostwald's Klassiker, no. 153, Leipzig, 1905.

{[}12{]} N. H. ABEL, Œuvres, 2 volumes, ed.~Sylow and Lie, Christiania, 1881.

{[}13{]} R. DEDEKIND, Gesammelte mathematische Werke, vol.~II, Braunschweig (Vieweg), 1932, p.~315.

{[}14{]} E. BOREL, Leçons sur la théorie des fonctions, 2nd edition, Paris (Gauthier-Villars), 1914.

{[}15{]} R. BAIRE, Leçons sur les fonctions discontinues, Paris (Gauthier-Villars), 1905.

{[}16{]} N. LUSIN, Leçons sur les ensembles analytiques et leurs applications, Paris (Gauthier-Villars), 1930.

\BourbakiBackSection{INDEX OF NOTATION}

The reference numbers indicate the chapter, section and sub-section (or, exercise) in that order.

\(\mathring{\mathrm{A}}, \overline{\mathrm{A}}\) (A a subset of a topological space): I, 1, 6.\\
\(\varprojlim \mathrm{X}_{\alpha}[(\mathrm{X}_{\alpha})\) an inverse system of topological spaces{]}: I, 4, 4.\\
\(\tilde{f}(a) (\tilde{f}\) the germ of a mapping): I, 6, 10.\\
\(\lim_{\mathfrak{G}} f, \lim_{x, \mathfrak{G}} f(x), \lim_{x} f(x): I, 7, 3.\) \(\lim_{n \to \infty} x_n: I, 7, 3.\) \(\lim_{x \in A} f(x)\) (A a directed set): I, 7, 3.\\
\(\lim_{x \to a} f(x): I, 7, 4.\) \(\lim_{x \to a, x \in A} f(x), \lim_{x \to a, x \neq a} f(x): I, 7, 5.\) \(f_{\mathrm{T}}: I, 5, 1.\) \(\iota_{\mathrm{X}}: I, 10.\)\\
Fr (A) (A a subset of a topological space): I, 1, Exercise 5.\\
\(\mathcal{C}_0(\mathrm{X}), \mathcal{C}_+(\mathrm{X}), \mathcal{C}_-(\mathrm{X}): I, 2,\) Exercise 5.\\
\(\mathfrak{P}_0(\mathrm{X}), \mathcal{C}_\Omega, \mathcal{C}_\Phi: I, 2,\) Exercise 7.\\
\(\mathfrak{F}(\mathrm{X}), \mathcal{C}_\Theta: I, 8,\) Exercise 12.\\
\(\mathfrak{C}^*: I, 8,\) Exercise 20.\\
\(\mathfrak{R}(\mathrm{X}): I, 9,\) Exercise 13.\\
\(\varprojlim \mathrm{X}_{\alpha}[(\mathrm{X}_{\alpha})\) an inverse system of uniform spaces{]}: II, 2, 7.\\
\(\hat{\mathrm{X}}\) (the Hausdorff completion of a uniform space X): II, 3, 7.\\
\(\tilde{\mathcal{U}}, \mathcal{C}(\tilde{\mathcal{U}}): II, 1,\) Exercise 5.\\
X/G (G a group operating on a space X): III, 2, 4.

P(K, L) (K, L subsets of a space with operators): III, 4, 5. \(\sum_{i \in I} x_i, \sum_{i} x_i, \sum x_i: III, 5, 1.\) \(\prod_{i \in I} x_i, \prod_{i} x_i, \prod x_i: III, 5, 1.\) \((x_n)\) (series, by abuse of language): III, 5, 6. \(\sum_{n=0}^{\infty} x_n, x_0 + x_1 + \cdots + x_n + \ldots, \sum_{n=0}^{\infty} x_n: III, 5, 6.\) \(\mathsf{P}_{n=0}^{\infty} x_n, \prod_{n=0}^{\infty} x_n: III, 5, 7.\) \(\mathbf{K}^*\) (K a division ring): III, 6, 7. \(\mathbf{R}, \mathbf{R}^*, \mathbf{R}_+, \mathbf{R}_*^+: IV, 1, 3.\) \(\sqrt[m]{x}, \sqrt{x}, x^{1/m}\) ( \(x\) real and \(\geqslant 0\) ): IV, 3, 3. \(\overline{\mathbf{R}}, +\infty, -\infty: IV, 4, 2.\) \(\sup A, \inf A\) (A a subset of \(\overline{\mathbf{R}}\) ): IV, 4, 2. \(\overline{\mathbf{R}}^*: IV, 4, 3.\) \(\lim_{x \to a, x > a, x \in A} f(x), \lim_{x \to a, x < a, x \in A} f(x): IV, 5, 3.\) \(f(a - ), f(a + ): IV, 5, 3.\) \(\sup_{x \in A} f(x), \inf_{x \in A} f(x)\) ( \(f\) a real-valued function): IV, 5, 4. \(\sup_{i \in I} f_i, \sup_{i} f_i, \inf_{i \in I} f_i, \inf_{i} f_i\) ( \(f_i\) real-valued functions): IV, 5, 5. \(\limsup_{\mathfrak{G}} f, \liminf_{\mathfrak{G}} f, \limsup_{x,\mathfrak{G}} f(x), \liminf_{x,\mathfrak{G}} f(x): IV, 5, 6.\) \(\limsup_{\mathfrak{G}} f, \liminf_{\mathfrak{G}} f, \limsup_{x} f(x), \liminf_{x} f(x): IV, 5, 6.\) \(\limsup_{x \to a} f(x), \liminf_{x \to a} f(x): IV, 5, 6.\) \(\limsup_{x \to a, x \in A} f(x), \liminf_{x \to a, x \in A} f(x), \limsup_{x \to a, x \neq a} f(x), \liminf_{x \to a, x \neq a} f(x): IV, 5, 6.\) \(\limsup_{n \to \infty} u_n, \liminf_{n \to \infty} u_n\) {[}(\(u_n\)) a sequence of numbers{]}: IV, 5, 6. \(\limsup_{n \to \infty} f_n, \liminf_{n \to \infty} f_n\) {[}(\(f_n\)) a sequence of functions{]}: IV, 5, 6. \(f + g, fg, 1/f\) ( \(f\) , \(g\) functions with values in \(\overline{\mathbf{R}}\) ): IV, 5, 7.

{[}x{]} (x a real number): IV, 8, 2.

\BourbakiBackSection{INDEX OF TERMINOLOGY}

The reference numbers indicate the chapter, section and sub-section (or exercise) in that order.

Absolute value of a real number : IV, 1, 6. Absolutely closed space : I, 9, Ex. 18. Absolutely convergent infinite product : IV, 7, 6. Absolutely convergent series : IV, 7, 6. Accessible space : I, 8, Ex. 1. Additive group of the rational line : IV, 1, 2. Additive group of the real line : IV, 1, 3. Additive uniformity of a topological division ring : III, 3, 8. Additive uniformity of the real line : IV, 1, 6. Alexandroff compactification : I, 9, 8. Alexandroff half-line : IV, 2, Exercise 12. Alexandroff's theorem : I, 9, 8. A-maximal topology : I, 3, Ex. 11. Algebraic complement of a subgroup : III, 6, 2. Alternating series : IV, 7, 6. Approximation to within ε by defect (excess) : IV, 8, 1. Arbitrarily small subgroups of a topological group : III, 2, Exercise 30. Archimedes' axiom : IV, 2, 1. Associated bijective continuous homomorphism : III, 2, 8. Associated Hausdorff group : III, 2, 6. Associated Hausdorff module : III, 6, 6. Associated Hausdorff ring : III, 6, 4. Associativity formula : III, 5, 3. Associativity of the sum of a summable family : III, 5, 3. Automorphism, local : III, 1, 3. Automorphism of a topological group : III, 1, 3. Axiom of Archimedes : IV, 2, 1. Axiom of Borel-Lebesgue : I, 9, 1. Axiom of Hausdorff : I, 8, 1.

\begin{Shaded}
\begin{Highlighting}[]
\NormalTok{Base of a filter : I, 6, 3.}
\NormalTok{Base (number) of an expansion of a real number : IV, 8, 5.}
\NormalTok{Base of a topology : I, 1, 3.}
\NormalTok{Base sequence of an expansion of a real number : IV, 8, 2.}
\NormalTok{Bicontinuous mapping : I, 2, 1.}
\NormalTok{Bolzano\textquotesingle{}s theorem : IV, 6, 1.}
\NormalTok{Borel{-}Lebesgue axiom : I, 9, 1.}
\NormalTok{Borel{-}Lebesgue theorem : IV, 2, 2.}
\NormalTok{Bounded above, below (function) : IV, 5, 1.}
\NormalTok{Bounded above, below (set) : IV, 2, 3.}
\NormalTok{Bounded function : IV, 5, 1.}
\NormalTok{Bounded (left, right) subset of a topological ring : III, 6, Exercise 12.}
\NormalTok{Bounded set (in a uniform space) : II, 4, Ex. 7.}
\NormalTok{Bounds, greatest lower and least upper (of a function) : IV, 5, 9.}

\NormalTok{Canonical mapping of the graph of the equivalence relation defined by a group operating freely : III, 4, 3.}
\NormalTok{Cantor\textquotesingle{}s theorem : IV, 8, 6.}
\NormalTok{Cantor\textquotesingle{}s triadic set : IV, 2, 5.}
\NormalTok{Cauchy filter : II, 3, 1.}
\NormalTok{Cauchy filter, minimal : II, 3, 2.}
\NormalTok{Cauchy sequence : II, 3, 1.}
\NormalTok{Cauchy\textquotesingle{}s condensation test : IV, 7, Exercise 3.}
\NormalTok{Cauchy\textquotesingle{}s criterion : II, 3, 3.}
\NormalTok{Cauchy\textquotesingle{}s criterion for series : III, 5, 46.}
\NormalTok{Cauchy\textquotesingle{}s criterion for summable families : III, 5, 2.}
\NormalTok{Chain (V{-}) : II, 4, 4.}
\NormalTok{Change of order of summation : III, 5, 3.}
\NormalTok{Close (V{-}) : II, 1, 1.}
\NormalTok{Closed covering : I, 1, 4.}
\NormalTok{Closed equivalence relation : I, 5, 2.}
\NormalTok{Closed mapping : I, 5, 1.}
\NormalTok{Closed set : I, 1, 4.}
\NormalTok{Closure of a set : I, 1, 6.}
\NormalTok{Cluster point of a filter base : I, 7, 2.}
\NormalTok{Cluster point of a function at a point relative to a subset : I, 7, 5.}
\NormalTok{Cluster point of a function with respect to a directed set : I, 7, 3.}
\NormalTok{Cluster point of a function with respect to a filter : I, 7, 3.}
\NormalTok{Cluster point of a function with respect to a sequence : I, 7, 3.}
\NormalTok{Cluster point of a germ of a function : I, 7, 3.}
\NormalTok{Coarser filter : I, 6, 2.}
\NormalTok{Coarser topology : I, 2, 2.}
\NormalTok{Coarser uniformity : II, 2, 2.}
\NormalTok{Coarsest topology : I, 2, 2.}
\end{Highlighting}
\end{Shaded}

Commutatively convergent series : III, 5, 7. Compact set : I, 9, 3. Compact space : I, 9, 3. Compactification, Alexandroff or one-point : I, 9, 8. Comparable filters : I, 6, 2. Comparable topologies : I, 2, 2. Comparable uniformities : II, 2, 2. Comparison principle for series : IV, 7, 1. Compatible (division ring structure and topology) : III, 6, 7. Compatible (equivalence relation and group of operators) : III, 2, 4. Compatible (group structure and topology) : III, 1, 1. Compatible (mapping of spaces with operators and homomorphism of groups of operators) : III, 2, 4. Compatible (ordered group structure and topology) : IV, 1, Exercise 1. Compatible (ring structure and topology) : III, 6, 3. Compatible (structure of group with operators and topology) : III, 6, 1. Compatible with a topology : II, 4, 1. Complement, algebraic : III, 6, 2. Complement, topological : III, 6, 2. Complete group : III, 3, 3. Complete ring : III, 6, 5. Complete uniform space : II, 3, 3. Completely Hausdorff (space, topology) : I, 9, Ex. 20. Completely irreducible between two points : IV, 2, Exercise 16. Completion of a Hausdorff topological division ring, field : III, 6, 8. Completion of a Hausdorff topological group : III, 3, 4. Completion of a Hausdorff topological module : III, 6, 6. Completion of a Hausdorff topological ring : III, 6, 5. Completion of a Hausdorff uniform space : II, 3, 7. Component, identity : III, 2, 2. Component of a point : I, 11, 5. Component of a subset : I, 11, 5. Condensation point : I, 9, Ex. 16. Condensation test, Cauchy's : IV, 7, Exercise 3. Connected component : I, 11, 5. Connected set : I, 11, 1. Connected space : I, 11, 1. Constituents, prime : II, 4, Ex. 19. Contiguous intervals : IV, 2, 5. Continuity, extension by : I, 8, 5. Continuous bijective homomorphism associated with a continuous homomorphism of topological groups : III, 2, 8. Continuous (mapping, function) : I, 2, 1. Continuous section : I, 3, 5.

Continuous with respect to a subspace : I, 3, 2. Continuously, group operating : III, 2, 4. Continuum, power of : IV, 8, 6. Convergent, absolutely : IV, 7, 6. Convergent, commutatively : III, 5, 7. Convergent filter : I, 7, 1. Convergent filter base : I, 7, 1. Convergent infinite product : III, 5, 7. Convergent infinite product of real numbers : IV, 7, 6. Convergent sequence : I, 7, 3. Convergent series : III, 5, 6. Convergent series of real numbers : IV, 7, 6. Coprefilter : I, 6, Ex. 17. Correspondence, proper : I, 10, Ex. 10. Covering, closed : I, 1, 4. Covering, open : I, 1, 1. Criterion, Cauchy's : II, 3, 3; III, 5, 2 and III, 5, 6. Cube root : IV, 3, 3.

Decimal expansion of a real number : IV, 8, 5. Dense set : I, 1, 6. Direct sum, topological : III, 6, 2. Dirichlet's function : IV, 6, 2. Discrete division ring : III, 6, 7. Discrete field : III, 6, 7. Discrete group : III, 1, 1. Discrete ring : III, 6, 3. Discrete (topological space, topology) : I, 1, 1. Discrete (uniform space, uniformity) : II, 1, 1. Division ring, topological : III, 6, 7. Dyadic expansion of a real number : IV, 8, 5.

Elementary filter : I, 6, 8. Elementary filter associated with a sequence : I, 6, 8. Elementary set : I, 4, 1. Ends of a locally compact space : I, 11, Ex. 19. Entourage : II, 1, 1. Entourage, symmetric : II, 1, 1. Entourages, fundamental system of : II, 1, 1. Envelope, lower (upper) : IV, 5, 5. Equivalence relation, closed : I, 5, 2. Equivalence relation defined by a group of operators : III, 2, 4. Equivalence relation, Hausdorff : I, 8, 3. Equivalence relation, open : I, 5, 2.

Expansion, decimal : IV, 8, 5.

Expansion, dyadic : IV, 8, 5.

Expansion, improper : IV, 8, 3.

Expansion of a real number with respect to a base sequence : IV, 8, 2.

Expansion, terminating : IV, 8, 3.

Expansion to base a : IV, 8, 5.

Expansion, triadic : IV, 8, 5.

Extended real line : IV, 4, 2.

Extension of a mapping by continuity : I, 8, 5.

Extension of identities, principle of : I, 8, 1.

Extension of inequalities, principle of : IV, 5, 2.

Exterior of a set : I, 1, 6.

Exterior point : I, 1, 6.

External semi-direct product : III, 2, 10.

External topological semi-direct product : III, 2, 10.

Extremally disconnected space : I, 11, Ex. 21.

Factor, general (of an infinite product): III, 5, 7.

Family, locally finite : I, 1, 5.

Family, multipliable : III, 5, 1.

Family, summable : III, 5, 1.

Field, discrete : III, 6, 7.

Field of real numbers : IV, 3, 1.

Field, p-adic : III, 6, Exercise 23.

Field, topological : III, 6, 7.

Filter : I, 6, 1.

Filter base : I, 6, 3.

Filter, Cauchy : II, 3, 1.

Filter coarser than another : I, 6, 2.

Filter comparable with another : I, 6, 2.

Filter, convergent : I, 7, 1.

Filter, elementary : I, 6, 8.

Filter, elementary, associated with a sequence : I, 6, 8.

Filter finer than another : I, 6, 2.

Filter, Fréchet : I, 6, 1.

Filter generated by a set of subsets : I, 6, 2.

Filter, induced : I, 6, 5.

Filter, intersection : I, 6, 2.

Filter, minimal Cauchy : II, 3, 2.

Filter, neighbourhood : I, 6, 1.

Filter, section : I, 6, 3.

Filter strictly coarser than another : I, 6, 2.

Filter strictly finer than another : I, 6, 2.

Filter subbase : I, 6, 2.

Filtered set : I, 6, 1.

Final topology : I, 2, 4.

Finer filter : I, 6, 2.

Finer topology : I, 2, 2.

Finer uniformity : II, 2, 2.

Finite open coverings, uniformity of : II, 4, 1.

Finite partial sum : III, 5, 1.

Finite partitions, uniformity of : II, 2, 2.

Finite real number : IV, 4, 2.

Finite real-valued function : IV, 5, 1.

Freely, group operating : III, 4, 3.

Fréchet filter : I, 6, 1.

Frontier of a set : I, 1, 6.

Frontier point : I, 1, 6.

Full subset : I, 11, Ex. 14.

Function, continuous : I, 2, 1.

Function, real-valued (or real): IV, 5, 1.

Function, uniformly continuous : II, 2, 1.

Fundamental system of entourages : II, 1, 1.

Fundamental system of neighbourhoods : I, 1, 3.

Gelfand ring : III, 6, Exercise 11.

General factor of an infinite product : III, 5, 7.

General term of a series : VI, 5, 6.

Geometric progression : IV, 7, 1.

Germ of a mapping at a point: I, 6, 10.

Germ of a mapping with respect to a filter : I, 6, 9.

Germ of a subset at a point : I, 6, 10.

Germ of a subset with respect to a filter : I, 6, 9.

Greatest lower bound of a real-valued function : IV, 5, 4.

Greatest lower bound of a set of topologies : I, 2, 4.

Greatest lower bound of a set of uniformities : II, 2, 5.

Group, complete : III, 3, 3.

Group, discrete : III, 1, 1.

Group, Hausdorff, associated with a topological group : III, 2, 6.

Group operating freely on a set : III, 4, 3.

Group operating trivially on a set : III, 2, 4.

Group, paratopological : III, 1, Exercise 4.

Group, product topological : III, 2, 9.

Group, quasi-topological : III, 2, Exercise 5.

Group, quotient topological : III, 2, 6.

Group, semi-topological : III, 1, Exercise 2.

Group, topological : III, 1, 1.

Group, topological, operating continuously on a topological space : III, 2, 4.

Group, topological, operating properly on a topological space: III, 4, 1. Group, topological, with operators: III, 6, 1.

Half-line, Alexandroff : IV, 2, Exercise 12. Hamel base : IV, 6, Exercise 1. Hausdorff completion of a topological group : III, 3, 4. Hausdorff completion of a topological module : III, 6, 6. Hausdorff completion of a topological ring : III, 6, 5. Hausdorff completion of a uniform space : II, 3, 7. Hausdorff equivalence relation: I, 8, 3. Hausdorff group associated with a topological group : III, 2, 6. Hausdorff module associated with a topological module : III, 6, 6. Hausdorff ring associated with a topological ring : III, 6, 4. Hausdorff space : I, 8, 1. Hausdorff topology : I, 8, 1. Hausdorff uniform space associated with a uniform space : II, 3, 8. Hausdorff uniformity : II, 1, 2. Hausdorff's axiom : I, 8, 1. Homeomorphic topological spaces : I, 1, 1. Homeomorphism : I, 1, 1. Homeomorphism, local : I, 11, Ex. 25. Homogeneous space : III, 2, 5.

Identification, space obtained by : I, 3, 4. Identities, principle of extension of : I, 8, 1. Identity component of a topological group : III, 2, 2. Image, inverse, of a topology : I, 2, 3. Image, inverse, of a uniformity : II, 2, 4. Improper expansion of a real number : IV, 8, 3. Induced by a uniformity (topology) : II, 1, 2. Induced filter : I, 6, 5. Induced topology : I, 2, 3. Induced uniformity : II, 2, 4. Infinite product, absolutely convergent : IV, 7, 6. Infinite product, convergent : III, 5, 7. Infinite product defined by the sequence ( \(x_n\) ) : III, 5, 7. Infinite product of real numbers, convergent : IV, 7, 6. Infinite product whose general factor is \(x_n\) : III, 5, 7. Initial topology : I, 2, 3. Initial uniformity : II, 2, 3. Inner automorphism of a topological group : III, 1, 3. Integer, p-adic : III, 6, Exercise 23. Integral part of a real number : IV, 8, 2. Interior of a set : I, 1, 6.

Interior point: I, 1, 6.

Intersection filter : I, 6, 2.

Intersection topology : I, 1, 6.

Inverse image of a topology : I, 2, 3.

Inverse image of a uniformity : II, 2, 4.

Inverse limit of an inverse system of groups (rings): III, 7, 1.

Inverse limit of an inverse system of sets endowed with internal (external) laws: III, 7, 1.

Inverse limit of an inverse system of topological groups (modules, rings): III, 7, 2.

Inverse limit of topological spaces : I, 4, 4.

Inverse limit of topologies : I, 4, 4.

Inverse limit of uniform spaces : II, 2, 7.

Inverse limit of uniformities : II, 2, 7.

Inverse system of groups (rings): III, 7, 1.

Inverse system of topological groups (rings): III, 7, 2.

Inverse system of topological spaces : I, 4, 4.

Inverse system of topologies : I, 4, 4.

Inverse system of uniform spaces : II, 2, 7.

Inverse system of uniformities : II, 2, 7.

Irrational numbers : IV, 1, 2.

Irreducible between two points : II, 4, Ex. 19.

Isodyne space : I, 2, Ex. 11.

Isolated point : I, 1, 6.

Isomorphic uniform spaces : II, 1, 1.

Isomorphism, local : III, 1, 3.

Isomorphism of a topological group onto a topological group : III, 1, 3.

Isomorphism of a uniform space onto another : II, 1, 1.

Kolmogoroff space : I, 1, Ex. 2.

Least upper bound of a real-valued function : IV, 5, 4.

Least upper bound of a set of topologies : I, 2, 3.

Least upper bound of a set of uniformities : II, 2, 5.

Left adherent point : IV, 2, Exercise 1.

Left bounded : III, 6, Exercise 12.

Left inverse of a linear mapping : III, 6, 2.

Left topology on an ordered set : I, 1, Ex. 2.

Left uniformity of a topological group : III, 3, 1.

Length of a bounded interval : IV, 1, 5.

Limit, inverse : III, 7, 1 and 2.

Limit, lower (upper): IV, 5, 6.

Limit on the left (right): IV, 5, 3.

Limit (point) of a filter : I, 7, 1.

Limit (point) of a filter base : I, 7, 1.

Limit (point) of a function at a point: I, 7, 5.

Limit (point) of a function relative to a subset: I, 7, 5.

Limit (point) of a function with respect to a directed set: I, 7, 3.

Limit (point) of a function with respect to a filter : I, 7, 3.

Limit (point) of a germ of a function : I, 7, 3.

Limit (point) of a sequence : I, 7, 3.

Lindelöf space : I, 9, Ex. 14.

Line, extended real : IV, 4, 2.

Line, rational : I, 1, 2; IV, 1, 2.

Line, real : IV, 1, 3.

Local automorphism of a topological group : III, 1, 3.

Local direct product of topological groups : III, 2, Exercise 26.

Local homeomorphism : I, 11, Ex. 25.

Local isomorphism of a topological group with another : III, 1, 3.

Local product of a family of topological groups, relative to a family of open normal subgroups : III, 2, Exercise 26.

Locally bounded topological ring : III, 6, Exercise 12.

Locally closed subspace : I, 3, 2.

Locally compact space : I, 9, 7.

Locally connected space : I, 11, 6.

Locally finite family : I, 1, 5.

Locally isomorphic topological groups : III, 1, 3.

Locally precompact Hausdorff topological group : III, 3, Exercise 8.

Locally quasi-compact space : I, 9, Ex. 29.

Locally retrobounded topology : III, 6, Exercise 22.

Lower envelope of a family of real-valued functions : IV, 5, 5.

Lower limit of a real-valued function with respect to a filter : IV, 5, 6.

Lower semi-continuous real-valued function : IV, 6, 2.

Lower semi-continuous regularization of a real-valued function : IV, 6, 2.

Mapping, bicontinuous : I, 2, 1.

Mapping, closed : I, 5, 1.

Mapping, continuous : I, 2, 1.

Mapping, open : I, 5, 1.

Mapping, proper : I, 10, 1.

Mapping, uniformly continuous : II, 2, 1.

Maximum, relative : IV, 6, 2.

Maximum, strict relative : IV, 5, Exercise 10.

Minimal Cauchy filter : II, 3, 2.

Minimal Hausdorff space : I, 9, Ex. 19.

Minimum, relative : IV, 6, 2.

Mittag-Leffler's theorem : II, 3, 5.

\begin{Shaded}
\begin{Highlighting}[]
\NormalTok{Module, topological : III, 6, 6.}
\NormalTok{Morphism of spaces with operators : III, 2, 4.}
\NormalTok{Morphism of topological groups : III, 2, 8.}
\NormalTok{Morphism, strict : III, 2, 8.}
\NormalTok{Monotone limit, theorem of the : IV, 5, 2.}
\NormalTok{Multipliable family of elements of a multiplicative group : III, 5, 1.}
\NormalTok{Multiplicative uniformities of a topological division ring : III, 6, 8.}

\NormalTok{Neighbourhood filter : I, 6, 1.}
\NormalTok{Neighbourhood of a point : I, 1, 2.}
\NormalTok{Neighbourhood of a set : I, 1, 2.}
\NormalTok{Neighbourhood, symmetric : III, 1, 2.}
\NormalTok{Neighbourhood (V{-}) : II, 1, 2.}
\NormalTok{Neighbourhoods, fundamental system of : I, 1, 3.}
\NormalTok{Nilpotent, topologically : III, 6, Exercise 9.}
\NormalTok{nth root : IV, 3, 3.}
\NormalTok{Number, finite real : IV, 5, 2.}
\NormalTok{Number, irrational : IV, 1, 3.}
\NormalTok{Number, p{-}adic : III, 6, Exercise 23.}
\NormalTok{Number, real : IV, 1, 3 and (by abuse of language) IV, 4, 2.}

\NormalTok{One{-}point compactification : I, 9, 8.}
\NormalTok{Open covering : I, 1, 1.}
\NormalTok{Open equivalence relation : I, 5, 2.}
\NormalTok{Open mapping : I, 5, 1.}
\NormalTok{Open set : I, 1, 1.}
\NormalTok{Opposite of a topological group : III, 1, 1.}
\NormalTok{Orbit space of a space with operators by its group of operators : III, 2, 4.}
\NormalTok{Oscillation of a real{-}valued function : IV, 6, Exercise 6.}

\NormalTok{p{-}adic field, integers, numbers, topology, valuation : III, 6, Exercise 23.}
\NormalTok{p{-}adic solenoid : III, 7, Exercise 6.}
\NormalTok{p{-}adic uniformity : II, 1, 1.}
\NormalTok{Paracompact space : I, 9, 10.}
\NormalTok{Paratopological group : III, 1, Exercise 4.}
\NormalTok{Partial sum : III, 5, 3.}
\NormalTok{Pasting together of sets : I, 2, 5.}
\NormalTok{Pasting together of topological spaces : I, 2, 5.}
\NormalTok{Perfect set : I, 1, 6.}
\NormalTok{Periodic expansion of a real number : IV, 8, Exercise 3.}
\NormalTok{Permutable (operations of two groups on a space) : III, 2, 4.}
\NormalTok{Poincaré{-}Volterra theorem : I, 11, 7.}
\NormalTok{Point : I, 1, 1.}
\NormalTok{Point, cluster : I, 7, 2.}
\end{Highlighting}
\end{Shaded}

Point, condensation : I, 9, Ex. 16. Point, exterior : I, 1, 6. Point, frontier : I, 1, 6. Point, interior : I, 1, 6. Point, isolated : I, 1, 6. Point, limit : I, 7, 1. Point, singular : II, 4, Ex. 19. Points, V-close : II, 1, 1. Power of the continuum : IV, 8, 6. Precompact set : II, 4, 2. Precompact space : II, 4, 2. Prefilter : I, 6, Ex. 17. Prime constituents : II, 4, Ex. 19. Prime constituents, space of : II, 4, Ex. 19. Prime prefilter : I, 6, Ex. 17. Primitive set : I, 7, Ex. 8. Principle of comparison of series : IV, 7, 1. Principle of extension of identities : I, 8, 1. Principle of extension of inequalities : IV, 5, 2. Product, external semi-direct : III, 2, 10. Product filter : I, 6, 7. Product, local (direct) : III, 2, Exercise 26. Product of a multipliable family : III, 5, 1. Product of topological groups : III, 2, 9. Product of topological groups with operators : III, 6, 1. Product of topological rings : III, 6, 4. Product, semi-direct : III, 2, 10. Product, topological semi-direct : III, 2, 10. Product topological space : I, 4, 1. Product topology : I, 4, 1. Product uniform space : II, 2, 6. Product uniformity : II, 2, 6. Proper correspondence : I, 10, Ex. 10. Proper mapping : I, 10, 1. Properly, group operating : III, 4, 1.

Quasi-compact set : I, 9, 3. Quasi-compact space : I, 9, 1. Quasi-maximal topology : I, 2, Ex. 6. Quasi-ring : III, 6, Exercise 20. Quasi-topological group : III, 2, Exercise 5. Quotient space : I, 3, 4. Quotient space of a space with operators by its group of operators : III, 2, 4. Quotient topological group : III, 2, 6.

Quotient topological group with operators : III, 6, 1.

Quotient topological ring : III, 6, 4.

Quotient topology : I, 3, 4.

Quotient topology of the topology of a group by a subgroup : III, 2, 5.

Quotient topology of the topology of a space with operators by its group of operators : III, 2, 4.

Radical of a commutative topological group : III, 2, Exercise 28.

Radical group : III, 2, Exercise 28.

Rational line : I, 1, 2; IV, 1, 2.

Rational number space of n dimensions : I, 4, 1.

Real function : IV, 5, 1.

Real line : IV, 1, 3.

Real line, extended : IV, 4, 2.

Real number : IV, 1, 3 and (by abuse of language) IV, 4, 2.

Real-valued function : IV, 5, 1.

Real-valued function bounded above (below): IV, 5, 4.

Real-valued function, finite : IV, 5, 1.

Real-valued function, lower (upper) semi-continuous : IV, 6, 2.

Regular open set : I, 8, Ex. 20.

Regular space : I, 8, 4.

Regular topology : I, 8, 4.

Regularization, lower (upper) semi-continuous : IV, 6, 2.

Relative maximum (minimum) : IV, 6, 2.

Relatively compact : I, 9, 3.

Relatively quasi-compact : I, 9, 3.

Residual closed subgroup : III, 2, Exercise 28.

Residue of a series : III, 5, 6.

Restricted associativity of series : III, 5, 7.

Retrobounded set : III, 6, Exercise 22.

Right bounded : III, 6, Exercise 12.

Right inverse of a linear mapping : III, 6, 2.

Right topology (on an ordered set) : I, 1, Ex. 2.

Right uniformity of a topological group : III, 3, 1.

Ring, complete : III, 6, 5.

Ring, discrete : III, 6, 3.

Ring, Gelfand : III, 6, Exercise 11.

Ring, Hausdorff, associated with a topological ring : III, 6, 4.

Ring, locally bounded topological : III, 6, Exercise 12.

Ring, product topological : III, 6, 4.

Ring, quotient topological : III, 6, 4.

Ring, topological : III, 6, 3.

Ring, topological, Hausdorff completion of : III, 6, 5.

Roots (square, cube, nth) : IV, 3, 3.

σ-compact : I, 9, 9.

Saturation of a subset with respect to a group to a group of operators : III, 2, 4.

Section, continuous : I, 3, 5.

Section filter : I, 6, 3.

Semi-continuous, lower (upper) : IV, 6, 2.

Semi-direct product of subgroups : III, 2, 10.

Semi-topological goup : III, 1, Exercise 2.

Semi-regular (space, topology) : I, 8, Ex. 20.

Sequence, Cauchy : II, 3, 1.

Sequence convergent : I, 7, 3.

Series, absolutely convergent : IV, 7, 6.

Series, alternating : IV, 7, 6.

Series, commutatively convergent : III, 5, 7.

Series, convergent : III, 5, 6 and IV, 7, 6.

Series defined by the sequence \((x_{n})\) : III, 5, 6.

Series whose general term is \(x_{n} :\) III, 5, 6.

Set, bounded : II, 4, Ex. 7.

Set bounded above (below): IV, 2, 3.

Set, Cantor's triadic : IV, 2, 5.

Set, closed : II, 1, 4.

Set, compact : I, 9, 3.

Set, connected : I, 11, 1.

Set, dense : I, 1, 6.

Set, elementary : I, 4, 1.

Set, filtered : I, 6, 1.

Set, locally closed : I, 3, 3.

Set, open : I, 1, 1.

Set, perfect : I, 1, 6.

Set, precompact : II, 4, 2.

Set, quasi-compact : I, 9, 3.

Set, regular open : I, 8, Ex. 20.

Set, relatively compact : I, 9, 3.

Set, relatively quasi-compact : I, 9, 3.

Set, totally disconnected : I, 11, 5.

Set underlying a topological space : I, 1, 1.

Sign of a real number : IV, 3, 2.

Singular point : II, 4, Ex. 19.

Small (V-): II, 3, 1.

Solenoid, p-adic : III, 7, Exercise 6.

Solvable space : I, 2, Ex. 11.

Space, absolutely closed : I, 9, Ex. 18.

Space, accessible : I, 8, Ex. i.

Space, compact : I, 9, 1.

Space, complete : II, 3, 3.

Space, completely Hausdorff: I, 9, Ex. 20.

Space, connected : I, 11, 1.

Space, extremely disconnected : I, 11, Ex. 21.

Space, Hausdorff : I, 8, 1.

Space, Hausdorff, associated with a uniform space : II, 3, 8.

Space, homogeneous : III, 2, 5.

Space, isodyne : I, 2, Ex. 11.

Space, Kolmogoroff: I, 1, Ex. 2.

Space, Lindelöf : I, 9, Ex. 14.

Space, locally compact : I, 9, 7.

Space, locally compact \(\sigma\) -compact : I, 9, 9.

Space, locally connected : I, 11, 6.

Space, locally quasi-compact : I, 9, Ex. 29.

Space, minimal Hausdorff: I, 9, Ex. 19.

Space, \(n\) -dimensional rational : I, 4, 1.

Space of orbits of a group operating continuously on a topological space : III, 2, 4.

Space of prime constituents : II, 4, Ex. 19.

Space, paracompact : I, 9, 10.

Space, precompact : II, 4, 2.

Space, product topological : I, 4, 1.

Space, product uniform : II, 2, 6.

Space, quasi-compact : I, 9, 1.

Space, quotient : I, 3, 4; III, 2, 4.

Space, rational n-dimensional : I, 4, 1.

Space, regular : I, 8, 4.

Space, semi-regular : I, 8, Ex. 20.

Space, solvable : I, 2, Ex. 11.

Space, submaximal : I, 8, Ex. 22.

Space, sum : I, 2, 4.

Space, topological : I, 1, 1.

Space, topological homogeneous : III, 2, 5.

Space, topological, underlying a uniform space : II, 1, 2.

Space, totally disconnected : I, 11, 5.

Space, ultrafilter : I, 9, Ex. 26.

Space, ultraregular : I, 8, Ex. 25.

Space, uniform : II, 1, 1.

Space, uniformizable : II, 4,1.

Square root : IV, 3, 3.

Stabilizer : III, 2, 4.

Strict morphism : III, 2, 8.

Strict relative maximum : IV, 5, Exercise 10.

Strictly coarser filter : I, 6, 2.

Strictly coarser topology : I, 2, 2.

Strictly coarser uniformity : II, 2, 2.

Strictly finer filter : I, 6, 2.

Strictly finer topology : I, 2, 2.

Strictly finer uniformity : II, 2, 2.

Subbase of a filter : I, 6, 2.

Subbase of a topology : I, 2, 3.

Submaximal (space, topology) : I, 8, Ex. 22.

Subspace (of a topological space): I, 3, 1.

Subspace, locally closed : I, 3, 3.

Subspace of a uniform space : II, 2, 4.

Sum, finite partial : III, 5, 1.

Sum of a family of points of a commutative topological group : III, 5, 1.

Sum of a series : III, 5, 6.

Sum, partial : III, 5, 3.

Sum (topological space, topology) : I, 2, 4.

Summable family : III, 5, 1.

Symmetric entourage : II, 1, 1.

Symmetric neighbourhood : III, 1, 2.

Symmetry : III, 1, 1.

Terminating expansion (of a real number): IV, 8, 3.

Theorem, Alexandroff's : I, 9, 8.

Theorem, Mittag-Leffler's : II, 3, 5.

Theorem, Poincaré-Volterra : I, 11, 7.

Theorem, Tychonoff's : I, 9, 5.

Theorem of Bolzano : IV, 6, 1.

Theorem of Borel-Lebesgue : IV, 2, 2.

Theorem of Cantor : IV, 8, 6.

Theorem of the monotone limit : IV, 5, 2.

Theorem of Weierstrass : IV, 6, 1.

Topological complement of a stable subgroup : III, 6, 2.

Topological direct sum of stable subgroups : III, 6, 2.

Topological division ring : III, 6, 7.

Topological field : III, 6, 7.

Topological group : III, 1, 1.

Topological group with operators : III, 6, 1.

Topological homogeneous space : III, 2, 5.

Topological module : III, 6, 6.

Topological ring : III, 6, 3.

Topological semi-direct product : III, 2, 10.

Topological space : I, 1, 1.

Topological structure : see Topology.

Topologically nilpotent (element, ideal) : III, 6, Exercise 9.

Topology : I, 1, 1. Topology, A-maximal : I, 3, Ex. 11. Topology associated with a filter : I, 6, 5. Topology coarser than another : I, 2, 2. Topology, coarsest : I, 2, 2. Topology comparable with another : I, 2, 2. Topology, completely Hausdorff : I, 9, Ex. 20. Topology, discrete : I, 1, 1. Topology, final : I, 2, 4. Topology finer than another : I, 2, 2. Topology generated by a set of subsets : I, 2, 3. Topology, Hausdorff : I, 8, 1. Topology, induced : I, 2, 3. Topology induced by a uniformity : II, 1, 2. Topology, initial : I, 2, 3. Topology, intersection : I, 1, 6. Topology, left (on an ordered set) : I, 1, Ex. 2. Topology, locally bounded : III, 6, Exercise 12. Topology, locally retrobounded : III, 6, Exercise 22. Topology, p-adic : III, 6, Exercise 23. Topology, product : I, 2, 3. Topology, quasi-maximal : I, 2, Ex. 6. Topology, quotient : I, 2, 4. Topology, regular : I, 8, 4. Topology, right (on an ordered set) : I, 1, Ex. 2. Topology, semi-regular, associated with a topology : I, 8, Ex. 20. Topology strictly coarser than another : I, 2, 2. Topology strictly finer than another : I, 2, 2. Topology, submaximal : I, 8, Ex. 22. Topology, sum : I, 2, 4. Topology, ultraregular : I, 8, Ex. 25. Topology, uniformizable : II, 4, 1. Totally adherent : I, 9, Ex. 17. Totally disconnected (set, space) : I, 11, 5. Triadic expansion of a real number : IV, 8, 5. Triadic set, Cantor's : IV, 2, 6. Triangle inequality : IV, 1, 1. Trivial ultrafilter : I, 6, 4. Trivially, group operating : III, 2, 4. Two-sided uniformity on a topological group : III, 2, Exercise 6. Tychonoff's theorem : I, 9, 5.

Ultrafilter : I, 6, 4

Ultrafilter space : I, 9, Ex. 26.

Ultrafilter, trivial : I, 6, 4.

Ultraregular (space, topology) : I, 8, Ex. 25.

Underlying set of a topological space : I, 1, 1.

Underlying topological space of a uniform space : II, 1, 2.

Uniform structure : see Uniformity.

Uniformities, multiplicative (of a topological division ring) : III, 6, 8.

Uniformity : II, 1, 1.

Uniformity, additive (of a topological division ring): III, 6, 8.

Uniformity, additive (of the real line) : IV, 1, 6.

Uniformity coarser than another : II, 2, 2.

Uniformity comparable with another : II, 2, 2.

Uniformity compatible with a topology : II, 4, 1.

Uniformity, discrete : II, 1, 1.

Uniformity finer than another : II, 2, 2.

Uniformity, Hausdorff : II, 1, 2.

Uniformity, induced : II, 2, 4.

Uniformity, left (right) (of a topological group): III, 3, 1.

Uniformity of finite open coverings : II, 4, 1.

Uniformity of finite partitions : II, 2, 2.

Uniformity, p-adic : II, 1, 1.

Uniformity, product : II, 2, 6.

Uniformity strictly coarser than another : II, 2, 2.

Uniformity strictly finer than another : II, 2, 2.

Uniformity, two-sided : III, 3, Exercise 6.

Uniformizable (topological space, topology) : II, 4, 1.

Uniformly bounded above (below) (family of real-valued functions): IV, 5, 5.

Uniformly continuous mapping : II, 2, 1.

Upper envelope of a family of real-valued functions : IV, 5, 5.

Upper limit of a real-valued function with respect to a filter : IV, 5, 5.

Upper semi-continuous real-valued function : IV, 6, 2.

Upper semi-continuous regularization of a real-valued function : IV, 6, 2.

Valuation, p-adic : III, 6, Exercise 23.

Value, absolute : IV, 1, 6.

Value of a germ of a mapping at a point: I, 6, 10.

V-chain : II, 4, 4.

V-close (points) : II, 1, 1.

V-neighbourhood of a set : II, 1, 2.

V-small : II, 3, 1.

Weierstrass's theorem : IV, 6, 1.
